\documentclass[a4paper, serif, 10pt]{moderncv}

%% moderncv
% color and style
\moderncvcolor{black}
\moderncvstyle{banking}

%% page layout/margins
\usepackage[top=2.5cm, bottom=2.5cm, left=1.5cm, right=1.5cm]{geometry}

\nopagenumbers{}

%% Babel config for Dutch language
% ref:
% - https://latex3.github.io/babel/guides/locale-dutch.html
\usepackage[dutch]{babel}

%% fontspec
\usepackage{fontspec}

\IfFontExistsTF{Linux Libertine}{
  \setmainfont[Ligatures={TeX, Common}, Numbers=OldStyle]{Linux Libertine}
}{}
\IfFontExistsTF{Linux Libertine O}{
  \setmainfont[Ligatures={TeX, Common}, Numbers=OldStyle]{Linux Libertine O}
}{}

%% CTeX
% CJK support, used to show CJK characters in the resume
%
% - heading=false: disable ctex heading
% - punct=kaiming: use kaiming punctuations styles for CJK
% - scheme=plain: use plain scheme, do not override \normalsize font size
% - space=auto: space settings for CJK characters
%
% CJK fonts are resolved per locale language, e.g. Noto Serif/Sans CJK JP
% for Japanese resumes, with fallback to the next region if not installed
%
% ref:
% - http://ctan.mirrorcatalogs.com/language/chinese/ctex/ctex.pdf
\usepackage[UTF8, heading=false, punct=kaiming, scheme=plain, space=auto]{ctex}

\IfFontExistsTF{Noto Serif CJK SC}{
  \setCJKmainfont{Noto Serif CJK SC}
}{}
\IfFontExistsTF{Noto Sans CJK SC}{
  \setCJKsansfont{Noto Sans CJK SC}
}{}

%% line spacing
\usepackage{setspace}
\setstretch{1.125}

%% URL styling - use normal text instead of monospace
\urlstyle{same}

%% Patch \httplink and \httpslink to handle full URLs
% The original moderncv macros always prepend http:// or https:// to URLs.
% This causes issues when the URL already has a protocol (e.g., https://example.com)
% resulting in malformed URLs like https://https://example.com.
% This patch checks if the URL already contains "://" and uses it directly if so.
% Note: We use \str_set:Nx (x-expansion) to expand macros like \@homepage first.
\ExplSyntaxOn
\str_new:N \g_yamlresume_url_str

\RenewDocumentCommand{\httpslink}{O{}m}{
  \str_set:Nx \g_yamlresume_url_str {#2}
  \str_if_in:NnTF \g_yamlresume_url_str {://}
    { \url{#2} }
    { \url{https://#2} }
}

\RenewDocumentCommand{\httplink}{O{}m}{
  \str_set:Nx \g_yamlresume_url_str {#2}
  \str_if_in:NnTF \g_yamlresume_url_str {://}
    { \url{#2} }
    { \url{http://#2} }
}
\ExplSyntaxOff

\name{Sanne de Vries}{}
\title{Data Engineer met focus op schaalbare dataplatformen}
\phone[mobile]{+31 6 12345678}
\email{sanne.devries@example.com}
\homepage{https://sanne-devries.nl}

\address{Keizersgracht 123, Amsterdam, Noord-Holland, Nederland, 1015 CJ}{}{}

\extrainfo{{\small \faLinkedin}\ \href{https://www.linkedin.com/in/sannedevries}{@sannedevries} {} {} {} • {} {} {} 
{\small \faGithub}\ \href{https://github.com/sannedevries}{@sannedevries}}

\begin{document}

\maketitle

\section{Basis Informatie}

\cvline{}{Data Engineer met 6+ jaar ervaring in het ontwerpen en bouwen van betrouwbare, schaalbare dataplatformen. Sterk in het vertalen van bedrijfsbehoeften naar robuuste ETL/ELT-pipelines en datamodellen.
%
\begin{itemize}
\item Ervaren met cloud-native oplossingen (AWS, Azure) en moderne data stacks zoals Spark, dbt en Airflow
\item Focus op datakwaliteit, governance en het leveren van betrouwbare analyses voor BI- en ML-teams
\item Teamspeler die graag kennis deelt en junior collega's begeleidt
\end{itemize}}
\section{Opleidingen}

\cventry{sep 2015–jul 2017}
        {Master, Data Science en Technologie, Score: 8.2}
        {Technische Universiteit Delft}
        {\url{https://www.tudelft.nl}}
        {}
        {\begin{itemize}
\item Specialisatie in grootschalige dataverwerking en datamodellering
\item Afstudeeronderzoek naar schaalbare streaming-pipelines voor sensordata
\item Ervaring opgedaan met Python, Spark, Kafka en cloudplatformen
\end{itemize}
\textbf{Cursussen}: Gedistribueerde systemen, Databasetechnologie, Machine Learning, Big Data Architectuur, Cloud Computing, Data-ethiek en Privacy}
\section{Werkervaring}

\cventry{jan 2021–Heden}
        {Senior Data Engineer}
        {Rabobank}
        {\url{https://www.rabobank.nl}}
        {}
        {\begin{itemize}
\item Ontwierp en implementeerde een cloud-native dataplatform op AWS voor de afdeling Risk \& Compliance
\item Bouwde schaalbare ETL-pipelines met Python, Spark en Airflow, goed voor circa 2 TB dagelijkse data
\item Introduceerde dbt voor datamodellering en verbeterde de datakwaliteit met geautomatiseerde tests
\item Werkt samen met data scientists en BI-analisten om dataproducten van idee tot productie te brengen
\item Bewaakt governance en privacy (AVG) en begeleidt junior data engineers
\end{itemize}
\textbf{Trefwoorden}: AWS, Spark, Airflow, dbt, Data Governance}

\cventry{mrt 2018–dec 2020}
        {Data Engineer}
        {KPN}
        {\url{https://www.kpn.com}}
        {}
        {\begin{itemize}
\item Ontwikkelde batch- en streaming-pipelines met Kafka, Spark en Python voor netwerk- en klantdata
\item Migreerde een on-premise datawarehouse naar Azure Data Lake en Synapse
\item Automatiseerde datakwaliteitscontroles en monitoring, waardoor incidenten met 30\% daalden
\item Werkte in een agile team en droeg bij aan architectuurkeuzes en code reviews
\end{itemize}
\textbf{Trefwoorden}: Kafka, Azure, Python, Streaming, Datakwaliteit}

\cventry{sep 2017–feb 2018}
        {Junior Data Analist}
        {Achmea}
        {\url{https://www.achmea.nl}}
        {}
        {\begin{itemize}
\item Bouwde rapportages en dashboards in Power BI voor schade- en klantdata
\item Schreef SQL-queries en automatiseerde rapportages met Python
\item Ondersteunde data-engineering bij het opschonen en modelleren van brondata
\end{itemize}
\textbf{Trefwoorden}: SQL, Power BI, Python, Rapportage}
\section{Talen}

\cvline{Nederlands}{Moedertaal of tweetalige vaardigheid}
\cvline{Engels}{Volledige professionele vaardigheid \hfill \textbf{Trefwoorden}: Cambridge C1, Werktaal}
\cvline{Duits}{Beperkte werkvaardigheid \hfill \textbf{Trefwoorden}: Goethe B1, Leesvaardig}
\section{Vaardigheden}

\cvline{Data Engineering}{Expert \hfill \textbf{Trefwoorden}: Python, SQL, Spark, Airflow, dbt}
\cvline{Cloudplatformen}{Geavanceerd \hfill \textbf{Trefwoorden}: AWS, Azure, Databricks, Snowflake}
\cvline{Datawarehousing}{Geavanceerd \hfill \textbf{Trefwoorden}: Redshift, Synapse, BigQuery, Kimball}
\cvline{Streaming en messaging}{Geavanceerd \hfill \textbf{Trefwoorden}: Kafka, Kinesis, Event Hubs}
\cvline{Datagovernance}{Intermediair \hfill \textbf{Trefwoorden}: Datakwaliteit, Lineage, AVG, Metadata}
\cvline{Programmeren}{Geavanceerd \hfill \textbf{Trefwoorden}: Python, Scala, Bash, Git}
\section{Onderscheidingen}

\cventry{dec 2023}
        {Rabobank}
        {Data Engineer of the Year}
        {}
        {}
        {Erkend voor het ontwerpen van een schaalbaar en kostenefficiënt dataplatform dat meerdere teams binnen Risk \& Compliance ondersteunt.}
\section{Certificaten}

\cventry{jun 2022}
        {Amazon Web Services}
        {AWS Certified Data Analytics - Specialty}
        {\url{https://aws.amazon.com/certification/}}
        {}
        {}

\cventry{nov 2020}
        {Microsoft}
        {Microsoft Certified: Azure Data Engineer Associate}
        {\url{https://learn.microsoft.com/en-us/certifications/azure-data-engineer/}}
        {}
        {}
\section{Publicaties}

\cventry{mrt 2022}
        {Schaalbare datapipelines in de cloud}
        {Data Magazine NL}
        {\url{https://www.datamagazine.nl/artikelen/schaalbare-datapipelines}}
        {}
        {\begin{itemize}
\item Praktijkartikel over het ontwerpen van betrouwbare en schaalbare datapipelines
\item Gaat in op orchestratie, datakwaliteit en kostenbeheersing in cloudomgevingen
\item Gebaseerd op ervaringen met AWS, Spark, Airflow en dbt
\end{itemize}}
\section{Referenties}

\cventry{\emaillink[m.jansen@tudelft.nl]{m.jansen@tudelft.nl}}
        {Hoogleraar Data Science}
        {Dr. Marleen Jansen}
        {+31 15 278 1234}
        {}
        {Sanne combineert sterke technische vaardigheden met een scherp oog voor datakwaliteit en governance. Zij is een betrouwbare teamspeler die complexe dataproblemen gestructureerd aanpakt.}
\section{Projecten}

\cventry{jan 2023–jun 2023}
        {Open-source dashboard voor het inzichtelijk maken van energieverbruik in Nederlandse wijken}
        {EnergieDashboard}
        {\url{https://github.com/sannedevries/energiedashboard}}
        {}
        {\begin{itemize}
\item Verzamelt en verwerkt openbare energie- en weerdata via API's
\item Pipeline gebouwd met Python, dbt en DuckDB, georkestreerd met Airflow
\item Dashboard in Streamlit voor bewoners en beleidsmakers
\end{itemize}
\textbf{Trefwoorden}: Open Data, Duurzaamheid, Python, Streamlit}
\section{Interesses}

\cvline{Duurzaamheid}{Energie, Circulaire economie, Open data}
\cvline{Sport}{Hardlopen, Wielrennen, Klimmen}
\cvline{Muziek}{Gitaar, Festivals}
\section{Vrijwilligerswerk}

\cventry{sep 2021–Heden}
        {Vrijwilliger Data Mentor}
        {Stichting Code voor Nederland}
        {\url{https://www.codevoornederland.nl}}
        {}
        {\begin{itemize}
\item Begeleidt jongeren en zij-instromers bij hun eerste stappen in data-analyse en Python
\item Organiseert maandelijkse workshops over datakwaliteit en open data
\item Droeg bij aan een open-source datapipeline voor het verwerken van gemeentelijke open data
\end{itemize}}
    
\end{document}